% fm-11-mdl-formal-foundations

\subsection{1. Formal Setting}

This chapter establishes the mathematical foundations for the Minimum Description Length (MDL) methodology employed throughout GOLDEN BRIDGE v27. All claims about $\varphi$-structured expressions and their relation to physical constants are grounded in the framework developed here. Notation is fixed for the remainder of the document.

\subsubsection{1.1 The Trinity Alphabet}

\textbf{Definition 1.1 (Symbolic alphabet \_\_M13\_\_).} The \textit{Trinity alphabet} is the set

\_\_M0\_\_

equipped with the following algebraic operations:

- \textbf{Ternary multiplication}: the standard integer product, closed on \_\_M14\_\_ (i.e., the product of any two elements lies in \_\_M15\_\_).
- \textbf{Addition modulo prime \_\_M16\_\_}: for a prime \_\_M17\_\_, addition is taken in \_\_M18\_\_ and then projected back to the balanced ternary representative in \_\_M19\_\_ by the map \_\_M20\_\_.
- \textbf{The \_\_M21\_\_-power generator}: the family of scalings \_\_M22\_\_, where \_\_M23\_\_ is the golden ratio. For \_\_M24\_\_, the element \_\_M25\_\_ is an \textit{external} generator acting on expressions by scalar multiplication; it does not itself belong to \_\_M26\_\_ but extends the alphabet to the augmented set \_\_M27\_\_.

\textit{Remark 1.1.} The canonical identity \_\_M28\_\_ holds exactly in \_\_M29\_\_ and serves as the structural anchor linking the ternary base to the $\varphi$-power generator. This identity is used throughout without further citation.

\subsubsection{1.2 The $\varphi$-Grammar}

\textbf{Definition 1.2 (Symbolic grammar \_\_M30\_\_).} The \textit{$\varphi$-grammar} \_\_M31\_\_ is a weighted context-free grammar (WCFG) over the terminal vocabulary

\_\_M1\_\_

where \_\_M32\_\_ denotes the \_\_M33\_\_-th Lucas number and \_\_M34\_\_ denotes the \_\_M35\_\_-th Fibonacci number. \_\_M36\_\_ satisfies the following conditions:

1. \textbf{Finite production set}: \_\_M37\_\_ has a finite set of productions \_\_M38\_\_, each of the form \_\_M39\_\_ where \_\_M40\_\_ is a non-terminal and \_\_M41\_\_ is a string over terminals and non-terminals.
2. \textbf{Bounded operator arity}: every operator symbol in \_\_M42\_\_ has arity at most \_\_M43\_\_ for some fixed \_\_M44\_\_ (in the current implementation, \_\_M45\_\_).
3. \textbf{Integer cost weights}: each production \_\_M46\_\_ carries an integer cost \_\_M47\_\_ representing the number of bits required to encode the application of that production under a fixed prefix-free code.
4. \textbf{Prefix-free encoding}: the cost assignment is consistent with Kraft's inequality, \_\_M48\_\_, ensuring that \_\_M49\_\_ defines a valid probability distribution over derivation trees.

\textit{Remark 1.2.} The grammar \_\_M50\_\_ is \textbf{pre-registered}: its productions and cost weights are fixed before any target observable is examined (see Section 5). Post-hoc modification of \_\_M51\_\_ in response to observed residuals voids any pre-registration claim.

\subsubsection{1.3 Description Length}

\textbf{Definition 1.3 (Grammar-relative description length \_\_M52\_\_).} For an expression \_\_M53\_\_ generated by \_\_M54\_\_, the \textit{description length} of \_\_M55\_\_ under \_\_M56\_\_ is

\_\_M2\_\_

where \_\_M57\_\_ is the set of all derivation trees of \_\_M58\_\_ under \_\_M59\_\_ and the sum runs over all productions used in a given derivation tree \_\_M60\_\_.

\textit{Remark 1.3 (Grammar-dependence).} \_\_M61\_\_ depends explicitly on \_\_M62\_\_. It is \textbf{not} Kolmogorov-invariant in an absolute sense; it is an upper bound on the true Kolmogorov complexity of \_\_M63\_\_, subject to the invariance theorem (Theorem 2.1 below) only up to an additive constant that depends on the complexity of \_\_M64\_\_ itself. We write \_\_M65\_\_ to emphasise this distinction, and elaborate in Section 2.

\noindent\rule{\textwidth}{0.4pt}

\subsection{2. Kolmogorov Invariance and the Additive Constant}

\subsubsection{2.1 Classical Invariance Theorem}

\textbf{Theorem 2.1 (Invariance Theorem; Solomonoff 1964, Kolmogorov 1965, Chaitin 1969).} \textit{Let \_\_M66\_\_ and \_\_M67\_\_ be two universal Turing machines (equivalently, two universal description methods). Then there exists a constant \_\_M68\_\_, depending only on \_\_M69\_\_ and \_\_M70\_\_ and not on \_\_M71\_\_, such that for all finite strings \_\_M72\_\_:}

\_\_M3\_\_

\textit{Proof sketch.} Because \_\_M73\_\_ can simulate \_\_M74\_\_ by a fixed prefix program of length \_\_M75\_\_, any \_\_M76\_\_-description of \_\_M77\_\_ translates to a \_\_M78\_\_-description at most \_\_M79\_\_ bits longer. The bound is symmetric by exchanging roles.

\textit{Standard references for Theorem 2.1}: the statement appears in its modern form in (\href{https://homepages.cwi.nl/~paulv/papers/mdlindbayeskolmcompl.pdf}{Grünwald \& Vitányi et al., CWI Technical Report}) and is discussed in the context of policy-guided tree search by (\href{https://doi.org/10.3390/e28040434}{Li, 2026, Entropy 28(4):434}).

\textbf{Conditions of applicability.} Theorem 2.1 applies to universal Turing machines. The grammar \_\_M80\_\_ is \textbf{not} a universal Turing machine; it is a fixed finite-arity WCFG. Consequently Theorem 2.1 does not directly apply to \_\_M81\_\_, and the additive constant \_\_M82\_\_ in our setting must be interpreted as the description-length cost of encoding \_\_M83\_\_ itself within a reference universal machine.

\subsubsection{2.2 The Additive Constant in Practice}

\textbf{Corollary 2.1 (Upper-bound relationship).} For any expression \_\_M84\_\_ derivable from \_\_M85\_\_:

\_\_M4\_\_

where \_\_M86\_\_ is the Kolmogorov complexity of the grammar specification itself.

\textit{Critical caveat.} The constant \_\_M87\_\_ can be \textbf{arbitrarily large} relative to the MDL differences we observe in practice. For the GOLDEN BRIDGE application domain --- comparing $\varphi$-monomial expressions of description length \_\_M88\_\_--\_\_M89\_\_ bits against a dataset of CODATA physical constants --- the additive constant \_\_M90\_\_ is of the same order of magnitude as the description lengths being compared, and possibly larger. This means:

1. \textbf{We cannot claim} that a $\varphi$-expression with low \_\_M91\_\_ has low \textit{true} Kolmogorov complexity. It has low grammar-relative description length, which is a weaker and grammar-dependent statement.
2. \textbf{Comparisons between \_\_M92\_\_ and \_\_M93\_\_} (a competing grammar) are meaningful only when both grammars are pre-registered, because the constant \_\_M94\_\_ is not observable.
3. \textbf{The correct interpretation} of a "low-MDL-cost match" is: \textit{within the reference machine defined by \_\_M95\_\_, this expression is a short description of the target constant.} This is a statement about the pre-registered grammar, not about objective simplicity.

Formally:

\_\_M5\_\_

\_\_M96\_\_ is an \textbf{upper bound} on \_\_M97\_\_ under the reference machine defined by \_\_M98\_\_. All claims in this document about "description length" refer exclusively to \_\_M99\_\_.

\noindent\rule{\textwidth}{0.4pt}

\subsection{3. The Shaw--Cohan--Eisenstein--Toutanova (2025) Result}

\subsubsection{3.1 Statement and Construction}

(\href{https://arxiv.org/abs/2509.22445}{Shaw, Cohan, Eisenstein, Toutanova, arXiv:2509.22445}) establish a formal bridge between Kolmogorov complexity and practical deep learning. Their main result, paraphrased for our context, is:

\textbf{Theorem 3.1 (Asymptotic optimality of variational MDL; Shaw et al. 2025, informal).} \textit{A minimizer of an asymptotically optimal description length objective achieves optimal compression, for any dataset, up to an additive constant, in the limit as model resource bounds increase.}

More precisely, Shaw et al. construct a \textit{variational MDL objective} of the form

\_\_M6\_\_

with the model prior parameterised as an \textbf{adaptive Gaussian mixture}: \_\_M100\_\_, where the mixture components \_\_M101\_\_ are themselves learned. This construction is both \textbf{tractable} and \textbf{differentiable}, enabling gradient-based optimisation of the full two-part description length.

\textbf{Conditions of applicability for Theorem 3.1.} The asymptotic guarantee holds:
- \textit{in the limit} as the model family's resource bounds (parameter count, depth, etc.) grow without bound;
- for \textit{any fixed dataset} of finite size;
- \textit{up to an additive constant} that depends on the complexity of the optimisation algorithm itself.

For finite-resource models (including \_\_M102\_\_), the theorem provides a theoretical \textit{ceiling} on achievable compression, not a guarantee on any particular finite-grammar search.

\subsubsection{3.2 The Optimisation Gap}

Shaw et al. explicitly identify a critical limitation (\href{https://arxiv.org/abs/2509.22445}{arXiv:2509.22445}): \textbf{"standard optimizers fail to find such solutions from a random initialization, highlighting key optimization challenges."} This is not a peripheral remark; it is a central finding of the paper. Even when the globally optimal description-length solution exists in principle, gradient-based search from random initialisation fails to locate it reliably.

\textbf{Implication for GOLDEN BRIDGE.} Theorem 3.1 \textit{justifies} the MDL framework asymptotically --- it tells us that, in principle, minimising description length is the right objective. It does \textbf{not} guarantee that our finite-grammar \_\_M103\_\_ search algorithm recovers the universal Kolmogorov minimum, nor even a good approximation to it. The \_\_M104\_\_ grammar is used as a \textbf{tractable, pre-registered upper bound}: a constrained search space that limits complexity inflation and ensures falsifiability, rather than a claim to universality.

This distinction is not a weakness of the approach; it is its epistemic honesty. The pre-registration protocol (Section 5) and the multiple-testing correction (Section 6) are the mechanisms by which we maintain scientific integrity despite the impossibility of true Kolmogorov minimisation.

\noindent\rule{\textwidth}{0.4pt}

\subsection{4. Bayesian Equivalence and the Choice of Prior}

\subsubsection{4.1 The MDL / Bayes Correspondence}

The two-part MDL coding cost \_\_M105\_\_ admits a direct Bayesian interpretation. Minimising the total description length is equivalent to maximising the log-posterior under the prior \_\_M106\_\_:

\_\_M7\_\_

since \_\_M107\_\_ when the code is optimal. This equivalence is standard; see (\href{https://homepages.cwi.nl/~paulv/papers/mdlindbayeskolmcompl.pdf}{Grünwald \& Vitányi et al.}) for a comprehensive treatment.

\textbf{Definition 4.1 ($\varphi$-grammar prior).} The probability distribution over hypotheses \_\_M108\_\_ implied by \_\_M109\_\_ is:

\_\_M8\_\_

This is the \textit{Solomonoff-style prior restricted to \_\_M110\_\_}. It assigns higher probability to expressions with shorter \_\_M111\_\_-derivations.

\textbf{Epistemic commitment.} Definition 4.1 constitutes a \textbf{strong epistemic commitment}: by adopting \_\_M112\_\_ as the prior, the analyst asserts that $\varphi$-structured hypotheses (those expressible in \_\_M113\_\_ with short descriptions) are \textit{a priori} more probable than equally accurate hypotheses expressed in a different grammar. This commitment must be stated explicitly and pre-registered before data are examined, because the prior is not data-neutral.

\subsubsection{4.2 Alternative Priors}

Several competing priors are relevant comparators and SHOULD be reported alongside any \_\_M114\_\_ result.

\textbf{Speed Prior} (\href{https://doi.org/10.1007/3-540-45435-7\_15}{Schmidhuber, 2002, DOI: 10.1007/3-540-45435-7\_15}):

\_\_M9\_\_

where \_\_M115\_\_ is the program length and \_\_M116\_\_ is the runtime of \_\_M117\_\_ on a reference machine. The Speed Prior is a computable approximation to the Solomonoff prior that penalises not just description length but computational cost. It favours hypotheses that are both short and fast, making it appropriate for physical theories that must be computationally efficient. Unlike \_\_M118\_\_, it does not privilege $\varphi$-structured expressions \textit{a priori}.

\textbf{Levin Search prior} (\href{https://doi.org/10.3390/e28040434}{Li, 2026, Entropy 28(4):434}): a time-bounded variant of the Solomonoff prior with explicit resource accounting. Under this prior, the description length of \_\_M119\_\_ includes a term for the time required by Levin's universal search to find \_\_M120\_\_, making the prior sensitive to the tractability of the hypothesis class.

\textbf{AIC and BIC} as frequentist baselines:

\_\_M10\_\_

where \_\_M121\_\_ is the number of free parameters in \_\_M122\_\_, \_\_M123\_\_ is the sample size, and \_\_M124\_\_ is the maximised likelihood. AIC and BIC are \textbf{mandatory baseline comparators} for any reported \_\_M125\_\_ MDL gain. A $\varphi$-expression that achieves lower \_\_M126\_\_ than a competing expression but worse AIC or BIC (after accounting for the cost of \_\_M127\_\_ itself) has not demonstrated a genuine compression gain.

\subsubsection{4.3 Prior Sensitivity and Reporting}

Every result table in GOLDEN BRIDGE v27 that reports a "low-MDL-cost expression" SHALL include:

\begin{tabular}{|l|l|}
\hline
Metric \& Value \\
\hline
\_\_M128\_\_ \& (bits, pre-registered grammar) \\
AIC comparator \& (for the same target, best non-$\varphi$ model) \\
BIC comparator \& (same) \\
Speed Prior rank \& (position of \_\_M129\_\_ in Speed Prior ordering, if computable) \\
\hline
\end{tabular}
Omitting these comparators is a pre-registered violation of the reporting protocol.

\noindent\rule{\textwidth}{0.4pt}

\subsection{5. The Catalog42 Pre-Registration Protocol}

\subsubsection{5.1 Motivation}

The $\varphi$-grammar \_\_M130\_\_ can generate a very large number of candidate expressions for any given target observable. Without a binding pre-registration protocol, the analyst could --- even without dishonest intent --- examine many expressions, select the one with the lowest residual, and report it as a discovery. The Catalog42 protocol is designed to make this impossible by fixing all analytical choices before the test is run.

\subsubsection{5.2 Mandatory Fields}

Each Catalog42 entry is a timestamped, immutable record with the following \textbf{eight required fields}:

\textbf{Field 1 --- Hypothesis ID.} A string of the form \texttt{Catalog42-NNN} where \texttt{NNN} is a zero-padded integer. The ID is derived from the SHA-256 hash of the concatenation of Fields 2--5; it is immutable once registered.

\textbf{Field 2 --- Closed-form expression.} The expression \_\_M131\_\_ in full \_\_M132\_\_ notation, with every operator and operand explicitly listed. Abbreviated or informal notation is not permitted. Example: \_\_M133\_\_.

\textbf{Field 3 --- Pre-computed \_\_M134\_\_.} The description length \_\_M135\_\_ in bits, computed from the derivation tree of Field 2 under the pre-registered \_\_M136\_\_. This value MUST be computed before the residual against the target is known to the analyst.

\textbf{Field 4 --- Target observable.} The CODATA constant or experimental quantity being modelled, including:
- Reference value and uncertainty in the format \_\_M137\_\_ (CODATA 2022 or later);
- The specific CODATA dataset version used;
- The dimensional units.
Example: \_\_M138\_\_ (CODATA 2022, dimensionless; see \href{https://en.wikipedia.org/wiki/Fine-structure\_constant}{Wikipedia Fine-structure constant}).

\textbf{Field 5 --- Sanctioned random seeds.} If any stochastic component is used in the search (e.g., stochastic grammar expansion, Monte Carlo estimation), all random seeds SHALL be listed. The following seeds are \textbf{permanently forbidden}: \_\_M139\_\_ (see Appendix B). Use of a forbidden seed voids the entry.

\textbf{Field 6 --- Stopping rule.} The experiment terminates upon the \textbf{first} of:
- (a) Reaching the pre-declared sample budget \_\_M140\_\_ (number of candidate expressions evaluated); or
- (b) The falsification condition in Field 8 is triggered.
No extension of \_\_M141\_\_ after commencement is permitted.

\textbf{Field 7 --- Effect-size threshold.} The minimum residual reduction, in units of CODATA \_\_M142\_\_, that constitutes a "pass". The threshold SHALL be set to \_\_M143\_\_ (see Section 7). Sub-$\sigma$ residuals SHALL NOT be claimed as evidence of a genuine match; they are consistent with coincidence.

\textbf{Field 8 --- Falsification condition (numeric, before-the-fact).} The entry is \textbf{FALSIFIED} if a non-$\varphi$ alternative grammar \_\_M144\_\_ achieves strictly lower description length \_\_M145\_\_ on the same data, provided that \_\_M146\_\_ was itself pre-registered \textbf{before} the residual of Field 2 was observed. The falsification condition SHALL include a numeric threshold, e.g.: "If there exists \_\_M147\_\_ with \_\_M148\_\_ for pre-specified \_\_M149\_\_, then Catalog42-NNN is FALSIFIED."

\subsubsection{5.3 Zenodo Timestamping}

Every Catalog42 entry SHALL be deposited on Zenodo and assigned a DOI \textbf{before} the analyst runs the residual computation against the target observable. The Zenodo deposit timestamp is the binding pre-registration timestamp. Post-hoc deposits --- even by one second --- are inadmissible. The project-level Zenodo record (\href{https://doi.org/10.5281/zenodo.19227877}{DOI: 10.5281/zenodo.19227877}) serves as the parent record; individual Catalog42 entries are linked as versions or related records.

\noindent\rule{\textwidth}{0.4pt}

\subsection{6. Multiple-Testing Correction}

\subsubsection{6.1 The False-Discovery Problem}

The Catalog42 ledger is designed to grow. By the time GOLDEN BRIDGE v27 reaches publication, it may contain hundreds of $\varphi$-monomial candidates across dozens of target observables. The false-discovery probability under naive per-entry significance testing at level \_\_M150\_\_ grows as \_\_M151\_\_ where \_\_M152\_\_ is the number of tests; for \_\_M153\_\_ and \_\_M154\_\_, this exceeds 99\%.

(\href{https://sites.stat.columbia.edu/gelman/research/unpublished/forking.pdf}{Gelman \& Loken, 2013}) identify a related but subtler problem: even without explicit "fishing", the \textit{garden of forking paths} --- the implicit space of analytical choices made before data are examined --- inflates the false-positive rate beyond the nominal level. Every choice in the \_\_M155\_\_ specification (which operators to include, which cost weights to assign, which target observables to test) represents a fork in this garden. The pre-registration protocol (Section 5) reduces, but does not eliminate, this inflation.

\subsubsection{6.2 Mandated Correction Procedure}

\textbf{Protocol 6.1 (BH-FDR correction).} The Benjamini--Hochberg procedure (\href{https://www.jstor.org/stable/2346101}{Benjamini \& Hochberg, 1995}) SHALL be applied across the \textbf{full Catalog42 cohort} at target FDR \_\_M156\_\_. Specifically:

1. Let \_\_M157\_\_ be the ordered \_\_M158\_\_-values from all \_\_M159\_\_ Catalog42 entries tested in a given cohort.
2. Let \_\_M160\_\_.
3. Reject (declare "pass") only entries \_\_M161\_\_.
4. All reported "low-MDL-cost matches" SHALL cite their FDR-adjusted \_\_M162\_\_-value alongside the raw \_\_M163\_\_-value.

\textbf{Protocol 6.2 (Cohort definition).} The cohort is the \textbf{entire Catalog42 ledger} at the time of analysis, not a per-target or per-observable subset. Cherry-picking a subset to reduce \_\_M164\_\_ and thereby lower the BH threshold is a pre-registered violation.

\subsubsection{6.3 Reporting Requirements}

Every result that claims a "low-MDL-cost match" for a CODATA constant SHALL report all three of the following:

1. \textbf{Cohort rank}: the rank of the match within the full Catalog42 cohort when entries are sorted by residual magnitude.
2. \textbf{FDR-adjusted \_\_M165\_\_-value}: from Protocol 6.1.
3. \textbf{Effect size in CODATA \_\_M166\_\_ units}: \_\_M167\_\_, where \_\_M168\_\_ is the value of the $\varphi$-expression and \_\_M169\_\_ is the CODATA uncertainty.

Omission of any of these three quantities is a pre-registered violation of the reporting protocol.

\noindent\rule{\textwidth}{0.4pt}

\subsection{7. Power Analysis and Sample-Size Justification}

\subsubsection{7.1 General Framework}

For each target observable, the following quantities SHALL be pre-registered in the corresponding Catalog42 entry (Field 7):

- \_\_M170\_\_: the current measurement uncertainty from CODATA 2022 (or the most recent available version at the time of entry creation).
- \_\_M171\_\_: the \textbf{minimum detectable residual gain}, expressed in absolute units. Any $\varphi$-expression whose residual exceeds \_\_M172\_\_ is \textbf{automatically REJECTED} regardless of its MDL score.
- \_\_M173\_\_: the maximum number of candidate expressions evaluated, pre-declared in Field 6.

\textbf{Definition 7.1 (Automatic rejection criterion).} Let \_\_M174\_\_ be the real-number value of a $\varphi$-expression and \_\_M175\_\_ the target observable. The expression is \textbf{automatically rejected} if

\_\_M11\_\_

Automatic rejection applies regardless of the MDL score, the FDR-adjusted \_\_M176\_\_-value, or any other metric.

\subsubsection{7.2 Example: The Fine-Structure Constant \_\_M177\_\_}

The CODATA 2022 value of the inverse fine-structure constant is (\href{https://en.wikipedia.org/wiki/Fine-structure\_constant}{Wikipedia, Fine-structure constant}):

\_\_M12\_\_

Thus:

- \_\_M178\_\_ (absolute),
- \_\_M179\_\_.

Any $\varphi$-expression \_\_M180\_\_ with \_\_M181\_\_ is automatically rejected. Note that Eddington's 1929 conjecture of \_\_M182\_\_ exactly (\href{https://en.wikipedia.org/wiki/Fine-structure\_constant}{Wikipedia, Fine-structure constant}) fails this criterion by a margin of approximately \_\_M183\_\_, roughly 1700 standard deviations. The historical lesson is clear: approximate numerical coincidences at coarse precision are not evidence of a structural relationship.

\textit{Note on the broader physics literature.} The scientific consensus, as reviewed in (\href{https://en.wikipedia.org/wiki/Fine-structure\_constant}{Wikipedia, Fine-structure constant}), is that "no numerological explanation has ever been accepted by the physics community." The $\varphi$-grammar framework is not a numerological claim; it is a pre-registered, falsifiable compression test with explicit error control. The distinction matters.

\subsubsection{7.3 Replication Standard}

A result that survives automatic rejection (Definition 7.1) and FDR correction (Protocol 6.1) is designated a \textbf{provenance evidence} candidate, not a confirmed result. Promotion from "provenance evidence" to "robust witness" requires:

- \textbf{Independent replication by \_\_M184\_\_ distinct silicon-design teams}, each reproducing the anchor witness from open RTL source (see Theorem 36.1, anchor byte 0x47C0 at \texttt{\{uio\\_out, uo\\_out\}} on reset) using independently compiled bitstreams.
- Each replicating team SHALL pre-register its own Catalog42-equivalent entry on Zenodo before running the test.
- Replication failures SHALL be reported in full.

\noindent\rule{\textwidth}{0.4pt}

\subsection{8. Blind Analysis Protocol}

\subsubsection{8.1 Rationale}

Blind analysis is standard practice in particle physics and clinical trials because the experimenter's knowledge of the result, even when honest, can influence downstream choices in ways that inflate the false-positive rate. The two-stage protocol below is adapted from particle-physics convention.

\subsubsection{8.2 Stage A: Development (Blind to Target)}

During Stage A, the analyst constructs and refines \_\_M185\_\_, assigns cost weights, and populates the Catalog42 ledger using \textbf{training constants} --- closed-form mathematical constants and CODATA quantities that are not the designated target observable.

\textbf{Permissible Stage-A activities:}
- Grammar design and cost-weight assignment.
- Testing expressions against training constants (e.g., \_\_M186\_\_, \_\_M187\_\_, \_\_M188\_\_, speed of light \_\_M189\_\_, Planck constant \_\_M190\_\_, etc.).
- Revising the grammar in response to training-set performance.
- Pre-registering Catalog42 entries for the designated target.

\textbf{Prohibited Stage-A activities:}
- Computing the residual of any Catalog42 entry against the designated target observable.
- Examining any dataset or measurement that includes the designated target.

The designated target (\_\_M191\_\_ in the running example) is \textbf{hidden from the analyst} during Stage A. If the target value is publicly known (as \_\_M192\_\_ is), the analyst SHALL designate a proxy blinder (a co-author not involved in grammar design) who holds the target value and computes residuals only in Stage B.

\subsubsection{8.3 Stage B: Unblinding}

Stage B commences \textbf{only after} the Catalog42 ledger is timestamped on Zenodo (Section 5.3). The sequence is:

1. Zenodo DOI is obtained and recorded.
2. The proxy blinder (or automated verification script) computes residuals for all pre-registered entries against the target.
3. Results are reported in full, including all rejections.
4. \textbf{Any post-hoc grammar modification} --- adding operators, adjusting cost weights, adding new Catalog42 entries --- \textbf{voids the pre-registration} for the target in question.

Stage B results are final. There is no provision for "Stage C" re-analysis of the same target under the same grammar version.

\noindent\rule{\textwidth}{0.4pt}

\subsection{9. What This Framework Does NOT Claim}

The following negative claims are \textbf{explicit and binding}. They are enumerated here to prevent misrepresentation in derived communications, press releases, or secondary literature.

1. \textbf{We do NOT claim} that $\varphi$ "explains" physical constants in any causal-mechanical sense. The fine-structure constant \_\_M193\_\_ arises from the quantum electrodynamic coupling between charged particles and photons; its value is determined by experiment, not by the structure of the golden ratio. A low-MDL-cost $\varphi$-expression is a \textit{compressed description}, not a \textit{derivation}.

2. \textbf{We do NOT claim} that \_\_M194\_\_ recovers the true Kolmogorov-minimal description of any physical constant. As established in Sections 2 and 3, \_\_M195\_\_ always, and the gap can be arbitrarily large. The grammar \_\_M196\_\_ is a pre-registered upper bound, not a universal compressor.

3. \textbf{We do NOT claim statistical significance} without FDR correction. Raw \_\_M197\_\_-values from individual Catalog42 entries are uninformative without the BH correction applied across the full cohort (Protocol 6.1). Any report that cites a \_\_M198\_\_-value without the FDR-adjusted \_\_M199\_\_-value is incomplete.

4. \textbf{We do NOT claim that residuals below CODATA \_\_M200\_\_ are physical evidence.} A $\varphi$-expression that matches \_\_M201\_\_ to within \_\_M202\_\_ is consistent with the hypothesis that the match is a coincidence arising from the density of $\varphi$-monomials in the neighbourhood of the CODATA value. Sub-$\sigma$ agreement is a necessary but not sufficient condition for any claim.

5. \textbf{We do NOT claim that the silicon anchor 0x47C0 "proves" anything} beyond a pre-registered design choice committed at a specific git hash. The anchor byte is a witness to the implementation of the $\varphi$-grammar in RTL; it is not evidence that $\varphi$ governs physical law. "Anchor byte witness" (the correct language) is strictly weaker than "proof of $\varphi$ in nature" (forbidden language).

6. \textbf{We do NOT claim} that the asymptotic guarantee of Theorem 3.1 (Shaw et al. 2025) applies to our finite-grammar, finite-data setting. The theorem applies in the limit of unbounded resources; for \_\_M203\_\_ with bounded arity \_\_M204\_\_ and finite production set, the theorem provides only asymptotic motivation.

7. \textbf{We do NOT claim} that the Speed Prior (Schmidhuber 2002) or the Levin Search prior (Li 2026) agree with \_\_M205\_\_. They are distinct priors with distinct philosophical commitments. The $\varphi$-grammar prior is one specific choice, and the results are prior-dependent.

\noindent\rule{\textwidth}{0.4pt}

\subsection{10. Concrete Checklist for Future Catalog42 Entries}

Every new Catalog42 entry MUST satisfy \textbf{all ten} items on the following checklist before it may appear in any GOLDEN BRIDGE publication, preprint, or public communication. The checklist SHALL be completed by the submitting author and countersigned by at least one co-author.

\begin{tabular}{|l|l|l|}
\hline
\section{\& Item \& Required Answer \\}
\hline
1 \& Has a Zenodo DOI been obtained for this entry \textbf{before} computing the target residual? \& Yes \\
2 \& Is the grammar \_\_M206\_\_ version (production set, cost weights, arity bound) fully specified and identical to the pre-registered version? \& Yes \\
3 \& Has FDR correction (BH at \_\_M207\_\_) been applied across the \textbf{full} Catalog42 cohort, not just this entry? \& Yes \\
4 \& Has an alternative-grammar comparator \_\_M208\_\_ been pre-registered and its description length \_\_M209\_\_ reported for the same target? \& Yes \\
5 \& Are AIC and BIC reported alongside \_\_M210\_\_ for the same target observable? \& Yes \\
6 \& Is the effect size reported in units of CODATA \_\_M211\_\_ (not relative units, not absolute units alone)? \& Yes \\
7 \& Has the blind-analysis protocol (Section 8) been followed, with Stage A and Stage B clearly separated and attested? \& Yes \\
8 \& Are the forbidden seeds \_\_M212\_\_ absent from the sanctioned seeds list (Field 5)? \& Yes \\
9 \& Is the stopping rule (Field 6) respected --- i.e., \_\_M213\_\_ was not extended after commencement? \& Yes \\
10 \& Is the entry's result reported in full, including all entries that failed automatic rejection (Definition 7.1), not just passing entries? \& Yes \\
\hline
\end{tabular}
\textit{A "No" on any item renders the entry inadmissible for publication.} The checklist SHALL be archived alongside the Catalog42 entry on Zenodo.

\noindent\rule{\textwidth}{0.4pt}

\subsection{References}

- Shaw, P., Cohan, J., Eisenstein, J., \& Toutanova, K. (2025). "Bridging Kolmogorov Complexity and Deep Learning: Asymptotically Optimal Description Length Objectives for Transformers." \textit{arXiv:2509.22445}. \href{https://doi.org/10.48550/arXiv.2509.22445}{https://doi.org/10.48550/arXiv.2509.22445}

- Grünwald, P., \& Vitányi, P. (n.d.). "Shannon Information and Kolmogorov Complexity." CWI Technical Report. \href{https://homepages.cwi.nl/~paulv/papers/mdlindbayeskolmcompl.pdf}{https://homepages.cwi.nl/~paulv/papers/mdlindbayeskolmcompl.pdf}

- Li, M. (2026). "From Levin's Universal Search to Policy-Guided Tree Search." \textit{Entropy} 28(4):434. \href{https://doi.org/10.3390/e28040434}{https://doi.org/10.3390/e28040434}

- Schmidhuber, J. (2002). "The Speed Prior: A New Simplicity Measure Yielding Near-Optimal Computable Predictions." In \textit{Proceedings of COLT 2002}, LNAI 2375. \href{https://doi.org/10.1007/3-540-45435-7\_15}{https://doi.org/10.1007/3-540-45435-7\_15}

- Benjamini, Y., \& Hochberg, Y. (1995). "Controlling the False Discovery Rate: A Practical and Powerful Approach to Multiple Testing." \textit{Journal of the Royal Statistical Society, Series B}, 57(1):289--300.

- Gelman, A., \& Loken, E. (2013). "The garden of forking paths: Why multiple comparisons can be a problem, even when there is no 'fishing expedition'." \href{https://sites.stat.columbia.edu/gelman/research/unpublished/forking.pdf}{https://sites.stat.columbia.edu/gelman/research/unpublished/forking.pdf}

- Wikipedia contributors. "Fine-structure constant." \textit{Wikipedia, The Free Encyclopedia}. \href{https://en.wikipedia.org/wiki/Fine-structure\_constant}{https://en.wikipedia.org/wiki/Fine-structure\_constant}

- GOLDEN BRIDGE v27 project record. Zenodo. \href{https://doi.org/10.5281/zenodo.19227877}{https://doi.org/10.5281/zenodo.19227877}
